This section describes the installation procedure for the webamc
package. This procedure has been realised on {\em Debian 13}.
%%%%%%%%%%%%%%%%%%%%
\subsection{Dependencies installation}
We will first update the system.
\begin{bash}
$ sudo apt update
$ sudo apt upgrade
\end{bash}

For a basic use of webamc you will need a python3 environment and
optionally the git system (if you wish to install webamc from the
development version).
\begin{bash}
$ sudo apt install python3 python3-pip python3-venv git
\end{bash}%$

If you install webamc on the web server side you will need the apache
web server, the postgresql database manager, and the wsgi module so
that our web server can execute our python scripts.
\begin{bash}
$ sudo apt install apache2 postgresql libapache2-mod-wsgi-py3 libpq-dev
\end{bash}%$

Now if you install webamc on the client side (i.e., to compile latex
MCQ that you intend to load to the web interface) you will need
a tex system able to compile tex files to pdf files:
\begin{bash}
$ sudo apt install texlive-latex-base texlive-extra-utils
\end{bash}%$
%%%%%%%%%%%%%%%%%%%%
\subsection{Webamc installation}
We will run python in a virtual environment \file{py3-venv} that we
create here in \file{/home/webamc}:
\begin{bash}
$ whoami 
webamc
$ cd
$ pwd
/home/webamc
$ python3 -m venv py3-venv
\end{bash}%$

We then activate the virtual environment.
\begin{bash}
$ source py3-venv/bin/activate
\end{bash}%$

We can finally install webamc either from the repository:
\begin{bash}
(py3-venv) $ git clone [*{\TOOLURLGIT}*]
(py3-venv) $ cd webamc
(py3-venv) $ pip3 install .
\end{bash}%$
or from the last available release:
\begin{bash}
(py3-venv) $ pip3 install [*{\TOOLURLRELEASE}*]
\end{bash}
